% GENERATED FILE. Edit canonical lessons, then run npm run generate:books.
% canonical-source-hash: fnv1a64:8c93bafb35c9ec62
% canonical-lessons: FA-W25-shin, FA-W25-sad, FA-W25-ta, FA-W25-za

\chapter{Letters That Share a Sound}
\label{ch:fa-letters-share-a-sound}
\hlchaptermodality{25}

\begin{chapteropening}
\textbf{By the end of this chapter:} \emph{I can write shin, sad, ta and za, and say which letter I already know sounds the same.}

Complete the last lesson of chapter 25: i can write shin, sad, ta and za, and say which letter I already know sounds the same.
\end{chapteropening}

\section[shin]{\fa{ش} (shin) — the letter that opens \fa{شما}}

\label{lesson:FA-W25-shin}



\begin{quote}
Before the new one: write \fa{ک} and \fa{گ} once each, and name them.

Now say \textbf{\fa{شما}} (\emph{shomâ}), \textquotedblleft{}you (respectful)\textquotedblright{}. You have read it for a long time. This lesson takes one letter out of it and puts it in your hand.
\end{quote}

\subsection*{Script you'll notice: \fa{ش}}
\textbf{\fa{ش}} — \emph{shin}, the sound \emph{sh}.

It opens \textbf{\fa{شما}} (\emph{shomâ}), the respectful \textquotedblleft{}you\textquotedblright{}. It is \textbf{\fa{س}}, which you already write, with \textbf{three dots above}.

\subsection*{Writing: \fa{ش}}
\hlblockfigure{\detokenize{figures/FA-W25-shin-filmstrip.pdf}}{How \fa{ش} is written, stroke by stroke}

\begin{itemize}
\raggedright
  \item \textbf{1.} start here: form the three teeth from right to left
  \item \textbf{2.} flow into the final bowl without lifting
  \item \textbf{3.} lift, then place the lower-left dot
  \item \textbf{4.} lift again and place the lower-right dot
  \item \textbf{5.} lift again and place the centered upper dot — and only now lift
\end{itemize}

\textbf{Pen lifts: 3.}

\begin{quote}
Stroke order is one attested teaching order; hands and teachers vary. Source: Persian Online, How to Write Persian Characters, \fa{ش} demonstration at 01:29–01:35 (Liberal Arts Instructional Technology Services, Univ. of Texas at Austin).
\end{quote}

\subsection*{Your turn}
\begin{itemize}
\raggedright
  \item \emph{Look:} at \textbf{\fa{شما}} and point at \fa{ش} inside it
  \item \emph{Trace it:} \fa{ش} three times, saying \emph{sh} as you finish each one
  \item \emph{Write it:} \fa{گ}, then \fa{ش}, from memory
  \item \emph{Read it:} \textbf{\fa{شما}} aloud once more
\end{itemize}

\begin{tcolorbox}[breakable,colback=teal!4,colframe=teal!35!black,title={Before you move on}]
Which letter is this — \fa{ش}? (\textbf{shin}, \emph{sh}.) Which word did it come from? (\textbf{\fa{شما}}, \emph{shomâ}, \textquotedblleft{}you (respectful)\textquotedblright{}.)
\end{tcolorbox}
\section[sâd]{\fa{ص} (sâd) — the letter that opens \fa{صد}}

\label{lesson:FA-W25-sad}



\begin{quote}
Before the new one: write \fa{گ} and \fa{ش} once each, and name them.

Now say \textbf{\fa{صد}} (\emph{sad}), \textquotedblleft{}a hundred\textquotedblright{}. You have read it for a long time. This lesson takes one letter out of it and puts it in your hand.
\end{quote}

\subsection*{Script you'll notice: \fa{ص}}
\textbf{\fa{ص}} — \emph{sâd}, the sound \emph{s}.

It opens \textbf{\fa{صد}} (\emph{sad}), \textquotedblleft{}a hundred\textquotedblright{}. In Persian it sounds \textbf{exactly like \fa{س}}: the two letters are spelled differently and said the same, because \textbf{\fa{ص}} came in with Arabic words and kept its Arabic spelling. The letter is a closed \textbf{oval} with a bowl after it.

\subsection*{Writing: \fa{ص}}
Put your pen on \fa{ص} and follow its line. Copy the shape you can see — slowly, and larger than it is printed.

\begin{quote}
This book does not yet tell you \textbf{where to start this letter or which way to travel}. It is not written down here until it can be written down with a source. Copying what is in front of you needs no such source, and it is how the shape gets into your hand in the meantime.
\end{quote}

\subsection*{Your turn}
\begin{itemize}
\raggedright
  \item \emph{Look:} at \textbf{\fa{صد}} and point at \fa{ص} inside it
  \item \emph{Trace it:} \fa{ص} three times, saying \emph{s} as you finish each one
  \item \emph{Write it:} \fa{ش}, then \fa{ص}, from memory
  \item \emph{Read it:} \textbf{\fa{صد}} aloud once more
\end{itemize}

\begin{tcolorbox}[breakable,colback=teal!4,colframe=teal!35!black,title={Before you move on}]
Which letter is this — \fa{ص}? (\textbf{sâd}, \emph{s}.) Which word did it come from? (\textbf{\fa{صد}}, \emph{sad}, \textquotedblleft{}a hundred\textquotedblright{}.)
\end{tcolorbox}
\section[tâ]{\fa{ط} (tâ) — the letter inside \fa{چطور}}

\label{lesson:FA-W25-ta}



\begin{quote}
Before the new one: write \fa{ش} and \fa{ص} once each, and name them.

Now say \textbf{\fa{چطور}} (\emph{chetor}), \textquotedblleft{}how\textquotedblright{}. You have read it for a long time. This lesson takes one letter out of it and puts it in your hand.
\end{quote}

\subsection*{Script you'll notice: \fa{ط}}
\textbf{\fa{ط}} — \emph{tâ}, the sound \emph{t}.

It is inside \textbf{\fa{چطور}} (\emph{chetor}), \textquotedblleft{}how\textquotedblright{}. Like \textbf{\fa{ص}}, it sounds \textbf{exactly like a letter you know}, here \textbf{\fa{ت}}. A closed oval on the line and a \textbf{tall upright} drawn from the top down.

\subsection*{Writing: \fa{ط}}
\hlblockfigure{\detokenize{figures/FA-W25-ta-filmstrip.pdf}}{How \fa{ط} is written, stroke by stroke}

\begin{itemize}
\raggedright
  \item \textbf{1.} start here: loop counterclockwise around the closed oval
  \item \textbf{2.} finish left along the baseline without lifting
  \item \textbf{3.} lift once, then draw the tall upright top-to-bottom — and only now lift
\end{itemize}

\textbf{Pen lifts: 1.}

\begin{quote}
Stroke order is one attested teaching order; hands and teachers vary. Source: Persian Online, How to Write Persian Characters, \fa{ط} demonstration at 01:54–01:56 (Liberal Arts Instructional Technology Services, Univ. of Texas at Austin).
\end{quote}

\subsection*{Your turn}
\begin{itemize}
\raggedright
  \item \emph{Look:} at \textbf{\fa{چطور}} and point at \fa{ط} inside it
  \item \emph{Trace it:} \fa{ط} three times, saying \emph{t} as you finish each one
  \item \emph{Write it:} \fa{ص}, then \fa{ط}, from memory
  \item \emph{Read it:} \textbf{\fa{چطور}} aloud once more
\end{itemize}

\begin{tcolorbox}[breakable,colback=teal!4,colframe=teal!35!black,title={Before you move on}]
Which letter is this — \fa{ط}? (\textbf{tâ}, \emph{t}.) Which word did it come from? (\textbf{\fa{چطور}}, \emph{chetor}, \textquotedblleft{}how\textquotedblright{}.)
\end{tcolorbox}
\section[zâ]{\fa{ظ} (zâ) — the letter that ends \fa{حافظ}}

\label{lesson:FA-W25-za}



\begin{quote}
Before the new one: write \fa{ص} and \fa{ط} once each, and name them.

Now say \textbf{\fa{حافظ}} (\emph{hâfez}), \textquotedblleft{}guardian\textquotedblright{}. You have read it for a long time. This lesson takes one letter out of it and puts it in your hand.
\end{quote}

\subsection*{Script you'll notice: \fa{ظ}}
\textbf{\fa{ظ}} — \emph{zâ}, the sound \emph{z}.

It ends \textbf{\fa{حافظ}} (\emph{hâfez}), the second half of \textbf{\fa{خدا} \fa{حافظ}}. It is \textbf{\fa{ط}} with \textbf{one dot above}, and in Persian it sounds exactly like \textbf{\fa{ز}}.

\subsection*{Writing: \fa{ظ}}
\hlblockfigure{\detokenize{figures/FA-W25-za-filmstrip.pdf}}{How \fa{ظ} is written, stroke by stroke}

\begin{itemize}
\raggedright
  \item \textbf{1.} start here: loop counterclockwise around the closed oval
  \item \textbf{2.} finish left along the baseline without lifting
  \item \textbf{3.} lift once, then draw the tall upright top-to-bottom
  \item \textbf{4.} lift again, then place the upper dot — and only now lift
\end{itemize}

\textbf{Pen lifts: 2.}

\begin{quote}
Stroke order is one attested teaching order; hands and teachers vary. Source: Persian Online, How to Write Persian Characters, \fa{ظ} demonstration at 01:57–01:59 (Liberal Arts Instructional Technology Services, Univ. of Texas at Austin).
\end{quote}

\subsection*{Your turn}
\begin{itemize}
\raggedright
  \item \emph{Look:} at \textbf{\fa{حافظ}} and point at \fa{ظ} inside it
  \item \emph{Trace it:} \fa{ظ} three times, saying \emph{z} as you finish each one
  \item \emph{Write it:} \fa{ط}, then \fa{ظ}, from memory
  \item \emph{Read it:} \textbf{\fa{حافظ}} aloud once more
\end{itemize}

\begin{tcolorbox}[breakable,colback=teal!4,colframe=teal!35!black,title={Before you move on}]
Which letter is this — \fa{ظ}? (\textbf{zâ}, \emph{z}.) Which word did it come from? (\textbf{\fa{حافظ}}, \emph{hâfez}, \textquotedblleft{}guardian\textquotedblright{}.)
\end{tcolorbox}
